The strongest current empirical result is that frozen ESM-2 embeddings improved predictive accuracy relative to handcrafted sequence features across all non-mean regressors and all evaluated split strategies. This supports the hypothesis that pretrained protein representations can be useful for enzyme optimal-pH prediction. The evidence is strongest for predictive performance, not for mechanistic interpretation.

The split-strategy results are more nuanced. The row-wise random split did contain exact duplicate leakage, yet the 90\% MMseqs2 cluster-safe split did not show worse MAE than the generated random split. Similarly, the organism-held-out split produced lower MAE than the random split. These findings do not support a simple statement that the controlled splits are always harder under the current protocol. The lower-identity MMseqs2 splits were also completed: ESM-2 histogram gradient boosting MAE increased from 0.6148 at 90\% identity to 0.6282 at 70\% and 0.6822 at 50\%. Interpretation remains limited because the organism-held-out test set has lower target-pH variance than the random test set, and cluster composition changes across identity thresholds.

The error analysis suggests that extreme pH labels are associated with large errors, but this pattern must be interpreted in light of label provenance. Zhou and Fang reported disputed pH values and protein-identity assignments in the same EpHod resource \cite{zhou2026caution}; for B5LXD7, the underlying Bispora study reports acidic activity measurements \cite{luo2009novel}, whereas the local file labels the record pH 11.5. This accession is among the largest errors in the 50\% cluster-safe and EC-held-out tests. We therefore do not automatically treat its disagreement as model failure. A leave-one-flagged-record-out sensitivity analysis is reported in \texttt{outputs/metrics/source\_label\_sensitivity.csv}; for the best ESM-2 models, excluding B5LXD7 changes MAE by only -0.0038 on the 50\% split and -0.0037 on the EC-held-out split. This analysis does not correct the label or independently remeasure the original experiment. More generally, extreme-pH error patterns remain compatible with regression toward the mean, but they are not evidence for a causal biochemical mechanism.

The lower-identity cluster experiments provide a more informative stress test. ESM-2 histogram gradient boosting achieved MAE 0.6282 on the 70\% cluster-safe split and 0.6822 on the 50\% split, compared with 0.6148 on the 90\% split. This degradation at lower sequence identity is consistent with a harder sequence-generalization problem, although it does not isolate sequence identity from differences in cluster composition and target distribution.

The EC-held-out experiment provides a complementary test of enzyme-class generalization. Its best ESM-2 model, random forest regression, achieved MAE 0.8321, compared with 0.6250 on the generated random split. This larger error indicates that holding out enzyme classes is considerably more difficult than the current sequence-cluster splits. The comparison should still be interpreted with care because EC groups vary in size and the EC-held-out target distribution may differ from the random test distribution.
